\chapter{How the design was reached}
This book is where the design lives until it is written in Lean. The owner of
the original system describes it as he built it. The description is written as
requirements and small models, and the LeanFM MCP server checks the models. Every
hole it finds is either a mistake in the transcription or a real choice in the
design. Corrections go back into the specification, and the questions nobody can
answer yet are kept in the last chapter. The source is written in one pass once
the security questions there are settled; the others shape features that can
ship later.

Each entry below says where it lives now. Once something is in a Lean file, that
file is the record and this book only explains it.

\section{Corrections that changed the specification}
Each entry was caught in conversation and corrected the design. Each says
where it lives: in a Lean file, or in this book only, until it is realized.
\begin{enumerate}\tightlist
\item \textbf{\texttt{visible\_dt} is local.} ``visibleDt means nothing to us. That
  was just a field attached to the record to determine the order in which the
  current shard will read it.'' It is not replicated; an import gets the importing
  shard's own value. (\texttt{Requirements.proto}, field 15.)
\item \textbf{It is assigned at commit.} ``That is when visibleDt is assigned, when
  it is committed and processed.'' Assigned at begin, the feed skips records, and
  the server shows how. (\texttt{Implementation.lean}, the writer.)
\item \textbf{Jobs run where the blob was created,} and everything they produce is
  a file: \texttt{plans.pdf-english.txt} (now \texttt{plans.pdf--english.txt},
  below). Imports never run jobs.
  (\texttt{Requirements.lean}, proof ``Jobs never run on imports''.)
\item \textbf{Policy returns read and write.} \texttt{Eval(jwt, rego)} decides
  ``whether they can even tell if the file exists''. A file the caller may not
  read is answered exactly as a missing one. (\texttt{Requirements.lean}.
  Superseded below: existence is governed by \texttt{list}, and only a caller
  who may not list a file gets the missing-file answer.)
\item \textbf{The window is concurrency, not distance.} 128 lanes means 128
  imports at once, finishing ``in some random order''; a failing file may sit far
  behind the cursor. (\texttt{Requirements.lean}, ``Imports are bounded''.)
\item \textbf{Imports are downloads.} A record from a peer is pulled with its blob
  by ranges, resumed by seeking, and exists only when the blob is complete.
  (\texttt{Requirements.lean}, task \texttt{import}.)
\item \textbf{Replication is ``how far up are you''.} ``You tell me how far up you
  are on replication, and I give you everything you have not seen.'' Marks per
  origin by \texttt{dt} became a cursor per peer in that peer's feed.
  (\texttt{Requirements.lean}; model \texttt{repl-cursor}.)
\item \textbf{Two clocks.} ``peer.visibleDt is not our.visibleDt.'' The cursor
  is kept in the peer's clock, in its own table, and moves like a Kafka consumer's
  committed offset: imports finishing out of order hold it still, and the one that
  completes the run moves it to the top of the finished, contiguous run.
  (\texttt{Implementation.lean}; model \texttt{lanes-contiguous}.)
\item \textbf{Policies are found by walking the tree.} A file's own
  \texttt{plans.doc--policy.rego}, else its folder's \texttt{.--policy.rego}, else
  the parent's, up to the root. Files have ordinary URLs under
  \texttt{/files/<org>/<user>}, and a user's home folder and its policy are made
  the first time their token is seen. (Superseded below: the walk builds a table
  rather than stopping at the nearest policy, and home folders belong to seats.
  \texttt{Requirements.lean} still names \texttt{plans.pdf.rego} and serves content
  by \texttt{id} and \texttt{dt}.)
\item \textbf{A folder's policy comes first, or its parent's.} ``If you write into
  an empty directory, supply a policy first! Or else, it gets its parent.'' Folders
  are made like \texttt{mkdir -p}: writing a file creates the missing folders on
  its path. (Book only.)
\item \textbf{Policy is a table built from the root down.} ``A hashtable that is
  populated from the root, with entries overwritten as you go down the tree to the
  file.'' Not the nearest policy alone: each level overwrites only the entries it
  sets, supplied ones like the banner and calculated ones like \texttt{read} and
  \texttt{write}. ``You might have access to a file, but not the policy about it'': the
  shard runs the walk with its own authority and the caller sees only the result.
  (Book only: \texttt{Requirements.lean} still has one \texttt{.rego} per file or
  folder.)
\item \textbf{Naming conventions instead of arguments.} ``Should directories be
  classified? Are directories implicitly created? We want sensible defaults, where
  you may put a classification over a directory or on an individual file, and use
  naming conventions to get out of these weird pickles.'' Folders are implicit; a
  folder or a file may carry a policy; \texttt{name--}\textit{x} is metadata about
  \texttt{name} with its permissions by default, and
  \texttt{name--policy.rego--policy.rego} restricts who sees the policy. (Book
  only.)
\item \textbf{Names are content.} From an example file whose name alone accused a
  named person of a crime: ``I wonder what is in it?'' A name is governed like
  the content, so no channel may carry a name its reader may not see, and the
  permission over names narrows down the path. (That permission became
  \texttt{list} in the next entry. Book only.)
\item \textbf{\texttt{list} is its own permission.} ``Another permission: List,
  which can be granted so that you can see that a file exists, and ask for
  permission. Different classification systems (medical, GDPR, foreign military)
  can be wildly different.'' And: ``It is one flexible classification system,
  such that the files protected are heterogeneous.'' The lesson behind it: a
  product tied to NATO markings had to be rewritten to demonstrate GDPR cases to
  an audience in Amsterdam. The name is governed by \texttt{list}, the content by
  \texttt{read}; the server interprets only a core of the table. (Book only.)
\item \textbf{Roles are seats.} ``Roles are like seats, into which butts are
  placed. Better to assign files to a role, that a person fills. That way, when
  people rotate out and move on, you don't need a full-time offboarding guy to
  make sure information is not lost.'' Asking for access is out of band; the org
  chart is implicit in how roles are created, and the token setup ensures that a
  boss or a security officer can still get to things. Home folders are seats too;
  token life is ``how immediately do people get fired?'', a few hours for some
  uses. (Book only.)
\item \textbf{Many issuers.} ``Imagine that you trust JWT CAs: USA, EU. Somebody
  else trusts CN, RU. Someone else trusts US, CA, 711.'' A token from a local store
  that checked a driver's licence can say ``over 21, Virginia'' and nothing more.
  ``You will not necessarily merge every shard. But you are certainly free to pick
  shards and JWT signers that you trust.'' And: ``The shards you have should be the
  same everywhere, given that the shard trusts your cert enough to use the system
  at all.'' (Book only.)
\item \textbf{The point: a CMS you use with \texttt{curl}.} ``Subject all uploaded
  files to an AI chain that makes derived files out of raw files. If you want to
  upload map layers, upload the viewer that goes around it. The idea is tons of
  React apps, tons of services. People assemble apps out of everything they have
  access to, and the app development can be done in the restricted environment,
  without necessarily needing a dev team that is cleared for anything.'' Search
  hits carry their context and their position, such as forty seconds into a
  video. And what the audience really thinks: ``upload a whole bunch of security
  cameras; let AIs put vtt on it''. (Book only.)
\item \textbf{Apps read, APIs write.} ``Usually app viewers are read-only to its
  users. APIs use the user's id to look up what their permissions actually are,
  not necessarily all in the JWT. The APIs need to be pretty airtight on
  security.'' (Book only.)
\item \textbf{A map, too.} An earlier map system: feeds of moving and fixed objects,
  news turned by an LLM into reports with links and things on the map, and a
  ``what is this?'' on every object. And its lessons: ``no permission system;
  10 s query latency; no way to tell people `no public LLMs to code'; but we did
  have a `no Chinese models' rule.'' Twitter feeds were a source too, and trying
  that API for a month cost \$2000 by surprise: an external feed needs a hard
  budget. A process that pretended to be an airport, feeding out track updates to
  make planes fly around, stress-tested the map. (Book only.)
\item \textbf{Cursors over API feeds.} ``Put stuff behind APIs, like /apis/, and
  /files/: the apps, and the data plus apps. Allow users to do some of this stuff
  themselves. Fast APIs meant supporting multiple users using cursors over data
  feeds coming out of the APIs. Instead, we had a giant Python app that was
  constantly updated to impress whoever it was we last had a meeting with.''
  (Book only.)
\item \textbf{The door forwards.} ``That is what the ingressHttps is for: it
  forwards requests to internal http, to mix and match the APIs in the backend
  with the mounted React apps and files in the frontend.'' Its global type: browser
  to door \texttt{GET /apis/\$x/\$task}, door to \texttt{host[\$x]} \texttt{GET
  /\$task}, like the reverse proxy in microCMS; with big streams and websockets.
  (Protocols in \texttt{Session.lean}: \texttt{Forward}, \texttt{Socket}; its
  negative requirements are book only.)
\end{enumerate}

\section{Decisions and their reasons}
\begin{longtable}{>{\raggedright}p{0.24\textwidth}>{\raggedright}p{0.44\textwidth}>{\raggedright\arraybackslash}p{0.22\textwidth}}
\toprule
Decision & Why & Lives in\\
\midrule
\endhead
Writes never wait for another island & Islands get cut apart for a day: ``somebody
always dragged an anchor across the Mediterranean.'' A file server that stops
taking uploads is worse than one that reconciles later. & Requirements; model
\texttt{repl-ordered} (writes commit while apart)\\
\addlinespace
Records are immutable: \texttt{(id, dt)} with \texttt{prev\_dts} & The original
MongoDB design. Concurrent versions are kept, never overwritten, so nothing is lost
in a merge. & Requirements; proto\\
\addlinespace
Latest is the greatest \texttt{(dt, shard)} & A function of the records held, so
every shard holding the same records shows the same file. & Requirements\\
\addlinespace
Mutable fields are local & ``Immutable with exceptions to keep the schema
simple'': one record type, with \texttt{visible\_dt} ignored when imported. &
Proto\\
\addlinespace
A cursor per peer, in the peer's feed & A late commit lands at the end of the
peer's feed, past every cursor already given out. No origin needs to commit in
\texttt{dt} order, and records from a third shard are relayed. & Requirements;
\texttt{repl-cursor} proved\\
\addlinespace
Delivery at least once & Exactly once needs coordination the islands cannot
afford. ``It's ok if we get a bunch of files that exist and are the right size; we
just have to cursor over them.'' & Requirements; Implementation\\
\addlinespace
Blob before record & ``Pull the blob in a task that makes the record exist.'' A
visible record always has its blob, so a download never waits on another island.
After an outage the network may be ``going crazy catching up'', so pulls are
bounded rather than greedy. & Requirements\\
\addlinespace
128 lanes, 32 kept for small files & Bounded concurrency handles ``incredible
bursts'', and a short file always opens a lane. & Requirements;
\texttt{sched-reserved} proved\\
\addlinespace
A failing import yields its lane & Otherwise ``it would effectively monopolize one
lane as long as it tries''. & Requirements; \texttt{retry-yield} proved\\
\addlinespace
Uploads in parts & Proxies have body-size limits: ``make a zero length file, then
keep appending to it with different http calls, then tell it you are done.'' &
Requirements; Implementation\\
\addlinespace
Ranges are mandatory & ``Even if you do not ask for a range, sometimes the browser
will offer a range.'' Seeking in a 9000 MB video was instant because of it. &
Requirements; Implementation (\texttt{ServeContent})\\
\addlinespace
Encryption is the volume's & The original wrote seekable ciphertext by hand:
``unimaginable pain; I do not need to do that again.'' With LUKS or the cloud
disk, a range stays a plain file read. & Implementation\\
\addlinespace
Pure-Go SQLite & The original's Postgres cgo driver ``made build times impossibly
long''. \texttt{modernc.org/sqlite} keeps the original indexes. &
Implementation\\
\addlinespace
Policy as \texttt{Eval(jwt, rego)} & No marking system is built in, so a
GDPR demo needs no rewrite of a NATO product. The banner is computed per caller:
\texttt{\{list, read, write, banner, bg, fg, rank\}}. The page shows the greatest rank
at top and bottom. & Requirements; proto\\
\addlinespace
Concurrent policies must all allow & With latest-wins, a later edit on the other
island undoes a tightening. & Requirements; \texttt{policy-allHeads} proved\\
\addlinespace
One session per page & A peer that must know when the puller stops is a protocol
bug; each page is answered on its own. & \texttt{Session.lean}\\
\bottomrule
\end{longtable}
